% GENERATED FILE. Edit canonical lessons, then run npm run generate:books.
% canonical-source-hash: fnv1a64:ed0f8050e669fb27
% canonical-lessons: BN-C134-bete, BN-C134-jhal, BN-C134-nonta, BN-C134-shokto, BN-C134-norom

\chapter{Short, Spicy, Salty, Hard, Soft}
\label{ch:bn-short-spicy-salty-hard-soft}
\hlchaptermodality{134}

\begin{chapteropening}
\textbf{By the end of this chapter:} \emph{I can say short, spicy, salty, hard and soft.}

Complete the last lesson of chapter 134: i can say short, spicy, salty, hard and soft.
\end{chapteropening}

\section[\texorpdfstring{\bn{বেঁটে} (beṭe)}{বেঁটে (beṭe)}]{\bn{বেঁটে} (beṭe) — short (of height)}

\label{lesson:BN-C134-bete}



\begin{quote}
Before the new one: say the Bengali for a vase, then the Bengali for a cup.
\end{quote}

\subsection*{You'll want to know: \bn{বেঁটে}}
\textbf{\bn{বেঁটে}} — \emph{beṭe} — \textquotedblleft{}short (of height)\textquotedblright{}.

\begin{grammarlens}[title={What you've built}]
One more word for this chapter.
\end{grammarlens}

\subsection*{Your turn}
\begin{itemize}
\raggedright
  \item \emph{Say it:} \emph{beṭe}
  \item \emph{Say it:} \emph{beṭe}, once more
  \item \emph{Say it:} say \emph{beṭe}
  \item \emph{Recall:} say the Bengali for a vase, then the Bengali for a cup, then say \emph{beṭe} again
\end{itemize}

\begin{tcolorbox}[breakable,colback=teal!4,colframe=teal!35!black,title={Before you move on}]
What does \bn{বেঁটে} mean? (\textquotedblleft{}Short (of height)\textquotedblright{}.) Say it once more.
\end{tcolorbox}
\section[\texorpdfstring{\bn{ঝাল} (jhāl)}{ঝাল (jhāl)}]{\bn{ঝাল} (jhāl) — spicy}

\label{lesson:BN-C134-jhal}



\begin{quote}
Before the new one: say the Bengali for a cup, then the Bengali for short.
\end{quote}

\subsection*{You'll want to know: \bn{ঝাল}}
\textbf{\bn{ঝাল}} — \emph{jhāl} — \textquotedblleft{}spicy\textquotedblright{}.

\begin{grammarlens}[title={What you've built}]
One more word for this chapter.
\end{grammarlens}

\subsection*{Your turn}
\begin{itemize}
\raggedright
  \item \emph{Say it:} \emph{jhāl}
  \item \emph{Say it:} \emph{jhāl}, once more
  \item \emph{Say it:} say \emph{jhāl}
  \item \emph{Recall:} say the Bengali for a cup, then the Bengali for short, then say \emph{jhāl} again
\end{itemize}

\begin{tcolorbox}[breakable,colback=teal!4,colframe=teal!35!black,title={Before you move on}]
What does \bn{ঝাল} mean? (\textquotedblleft{}Spicy\textquotedblright{}.) Say it once more.
\end{tcolorbox}
\section[\texorpdfstring{\bn{নোনতা} (nontā)}{নোনতা (nontā)}]{\bn{নোনতা} (nontā) — salty}

\label{lesson:BN-C134-nonta}



\begin{quote}
Before the new one: say the Bengali for short, then the Bengali for spicy.
\end{quote}

\subsection*{You'll want to know: \bn{নোনতা}}
\textbf{\bn{নোনতা}} — \emph{nontā} — \textquotedblleft{}salty\textquotedblright{}.

\begin{grammarlens}[title={What you've built}]
One more word for this chapter.
\end{grammarlens}

\subsection*{Your turn}
\begin{itemize}
\raggedright
  \item \emph{Say it:} \emph{nontā}
  \item \emph{Say it:} \emph{nontā}, once more
  \item \emph{Say it:} say \emph{nontā}
  \item \emph{Recall:} say the Bengali for short, then the Bengali for spicy, then say \emph{nontā} again
\end{itemize}

\begin{tcolorbox}[breakable,colback=teal!4,colframe=teal!35!black,title={Before you move on}]
What does \bn{নোনতা} mean? (\textquotedblleft{}Salty\textquotedblright{}.) Say it once more.
\end{tcolorbox}
\section[\texorpdfstring{\bn{শক্ত} (shôktô)}{শক্ত (shôktô)}]{\bn{শক্ত} (shôktô) — hard}

\label{lesson:BN-C134-shokto}



\begin{quote}
Before the new one: say the Bengali for spicy, then the Bengali for salty.
\end{quote}

\subsection*{You'll want to know: \bn{শক্ত}}
\textbf{\bn{শক্ত}} — \emph{shôktô} — \textquotedblleft{}hard\textquotedblright{}.

\begin{grammarlens}[title={What you've built}]
One more word for this chapter.
\end{grammarlens}

\subsection*{Your turn}
\begin{itemize}
\raggedright
  \item \emph{Say it:} \emph{shôktô}
  \item \emph{Say it:} \emph{shôktô}, once more
  \item \emph{Say it:} say \emph{shôktô}
  \item \emph{Recall:} say the Bengali for spicy, then the Bengali for salty, then say \emph{shôktô} again
\end{itemize}

\begin{tcolorbox}[breakable,colback=teal!4,colframe=teal!35!black,title={Before you move on}]
What does \bn{শক্ত} mean? (\textquotedblleft{}Hard\textquotedblright{}.) Say it once more.
\end{tcolorbox}
\section[\texorpdfstring{\bn{নরম} (nôrôm)}{নরম (nôrôm)}]{\bn{নরম} (nôrôm) — soft}

\label{lesson:BN-C134-norom}



\begin{quote}
Before the new one: say the Bengali for salty, then the Bengali for hard.
\end{quote}

\subsection*{You'll want to know: \bn{নরম}}
\textbf{\bn{নরম}} — \emph{nôrôm} — \textquotedblleft{}soft\textquotedblright{}.

\begin{grammarlens}[title={What you've built}]
One more word for this chapter.
\end{grammarlens}

\subsection*{Your turn}
\begin{itemize}
\raggedright
  \item \emph{Say it:} \emph{nôrôm}
  \item \emph{Say it:} \emph{nôrôm}, once more
  \item \emph{Say it:} say \emph{nôrôm}
  \item \emph{Recall:} say the Bengali for salty, then the Bengali for hard, then say \emph{nôrôm} again
\end{itemize}

\begin{tcolorbox}[breakable,colback=teal!4,colframe=teal!35!black,title={Before you move on}]
What does \bn{নরম} mean? (\textquotedblleft{}Soft\textquotedblright{}.) Say it once more.
\end{tcolorbox}
